\subsection{The saturated crystalline ordinary component and its actual points}
\label{subsec:sr-saturated-ordinary-crystalline}
We prove the local assertion (GE7), including the nearest-weight
case \(a=1\).
This concerns the local ring used in the generic rank-three
deformation problem.
The Gaussian argument has its independent local separator.

Let \(G_\ell=G_{\mathbb Q_\ell}\), let \(\eta\) be the
unramified local sign of \(K/\mathbb Q\), possibly trivial, and let \(\tau\) be the
finite character on inertia obtained from \(\psi_\ell\).
In the principal packet \(\tau^2=1\).
Put
\[
 H=\ker\eta\cap\ker\tau=G_k,\qquad
 \Psi|_{H\cap I_\ell}=\epsilon_\ell^a,\qquad a\ge1 .
 \tag{LC1}
\]
Here \(k/\mathbb Q_\ell\) is finite.
Choose \(s_0\in H\cap I_\ell\) with
\(\epsilon_\ell(s_0)\) of infinite order, and choose
\(\sigma\in H\) mapping to a topological generator of
\(H/(H\cap I_\ell)\).
These choices exist because both kernels are open and
the cyclotomic image of every open inertia subgroup is infinite.
Write \(h=(a,0,-a)\) and
\(p_i=\Psi(s_0)^{(1,0,-1)_i}\).
The nonzero constants \(p_i-p_j\) become units in every
finite characteristic-zero coefficient algebra.

Let \(\mathcal R\) be the actual complete Noetherian
local universal framed ring, with universal semidirect
representation \(r\).
For \(u\in H\cap I_\ell\), put
\[
 S_i(u)=r(u)-\Psi(u)^{(1,0,-1)_i}1,\qquad
 c_i=(p_i-p_{i+1})(p_i-p_{i+2}),
\]
where the indices in \(c_i\) are taken modulo three.
Let \(J_{\rm ord}\) be generated by the coefficients of
\[
 \operatorname{char}_{r(u)}(T)
       -\prod_i(T-\Psi(u)^{(1,0,-1)_i})
\]
and by the entries of
\(S_0(u_0)S_1(u_1)S_2(u_2)\), for all
\(u,u_i\in H\cap I_\ell\).
Define
\[
 \begin{split}
 T_i&=\operatorname{tr}
       \bigl(S_{i+1}(s_0)S_{i+2}(s_0)r(\sigma)\bigr),\\
 d_i&=c_{i+1}T_i-\epsilon_\ell(\sigma)c_iT_{i+1}
                       \quad(i=0,1),\\
 \Delta&=\ell\,d_0d_1,\qquad
 I_Y=(J_{\rm ord}:\Delta^\infty),\qquad
 I_{\rm cr}=\sqrt{I_Y}.
 \end{split}
 \tag{LC2}
\]
These are the actual \(\mathrm{ordIdeal},dd,\mathrm{del},IY,Icr\);
in particular saturation is performed in the universal ring
before evaluating a point.

We use the following standard local inputs:
crystalline representations form a category closed under
subobjects, quotients and tensor products; an unramified
finite-dimensional \(\mathbb Q_\ell\)-representation is crystalline;
and the Bloch--Kato fundamental exact sequence computes \(H^1_f\).
The dimension formula from that sequence and local
Euler--Poincare duality will be used explicitly below.
These are the local comparison and cohomology inputs of the
manuscript, rather than an automorphy assertion.\footnote{See
S.~Bloch and K.~Kato,
\href{https://math.stanford.edu/~conrad/BSDseminar/refs/BKTamagawa.pdf}
{\emph{\(L\)-functions and Tamagawa numbers of motives}},
Section 3, for the fundamental sequence and its finite local condition.}

\begin{srlemma}\label{lem:sr-saturated-artinian-lift}
Let \(R\) be Noetherian, \(J=(J:\Delta)\), and suppose
\(\ell\mid\Delta\).
For every \(R\)-point \(f:R\to B\), with \(B\) a finite local
\(\mathbb Q_\ell\)-algebra and \(J\subset\ker f\), there are
a finite local \(\mathbb Q_\ell\)-algebra \(B'\), a map
\(f':R\to B'\) killing \(J\), and a map \(q:B'\to B\)
such that
\[
 qf'=f,\qquad q(\operatorname{Ann}_{B'}f'(\Delta))=0 .
 \tag{LC3}
\]
All maps preserve the given coefficient scalars.
\end{srlemma}
\begin{proof}
Set \(A=(R/J)[1/\ell]\), and let \(\delta\) be the image of
\(\Delta\).
Saturation says multiplication by \(\delta\) is injective on
\(R/J\), hence on \(A\).
The point induces \(g:A\to B\).
Put \(\mathfrak m=g^{-1}\mathfrak m_B\).
The algebra \(A/\mathfrak m\) embeds in the finite extension
\(B/\mathfrak m_B\) of \(\mathbb Q_\ell\), so is a finite
field extension itself; thus \(\mathfrak m\) is maximal.
Choose \(r\) with
\(g(\mathfrak m^r)=0\), using nilpotence of \(\mathfrak m_B\).
Artin--Rees for the submodule \(\delta A\subset A\) gives
\(s\ge r\) such that
\[
              (\mathfrak m^s:\delta)\subset\mathfrak m^r.
\]
Indeed
\(\delta A\cap\mathfrak m^s\subset\delta\mathfrak m^r\),
and \(\delta\) can be canceled in \(A\).
Take \(B'=A/\mathfrak m^s\).
It is local and finite over \(\mathbb Q_\ell\), since
\(\mathfrak m\) is finitely generated and its residue field
is finite over \(\mathbb Q_\ell\).
The induced \(q\) exists since
\(\mathfrak m^s\subset\ker g\).
If \(\delta b=0\) in \(B'\), a representative of \(b\)
belongs to \((\mathfrak m^s:\delta)\), and hence maps
to zero under \(g\).
This proves (LC3).
\end{proof}

\begin{srlemma}\label{lem:sr-local-positive-annihilator}
Let \(B'\) be a finite local \(\mathbb Q_\ell\)-algebra,
and let \(\chi:H\to B'^\times\) be continuous with
\(\chi|_{H\cap I_\ell}=\epsilon_\ell^w\), \(w\ge1\).
Every continuous cocycle \(z\in Z^1(H,B'(\chi))\) can be
written
\[
                 z=z_f+z_{\rm exc},\qquad
 z_f\in Z^1_f(H,B'(\chi)),
 \tag{LC4}
\]
where \(z_{\rm exc}\) takes values in
\[
 W=\begin{cases}
 \operatorname{Ann}_{B'}
       (\chi(\sigma)-\epsilon_\ell(\sigma)),&w=1,\\
 0,&w\ge2.
 \end{cases}
\]
Here \(Z^1_f\) is the inverse image of \(H^1_f\) in \(Z^1\).
\end{srlemma}
\begin{proof}
Choose a coefficient field \(E\subset B'\), lifting its
finite residue field; it exists by separability in characteristic zero.
Regard \(B'\) and \(W\) as \(E\)-linear representations.
Dividing \(\chi\) by \(\epsilon_\ell^w\) gives an unramified
representation, so both are crystalline.
Every composition factor has the positive cyclotomic exponent \(w\).
Consequently \(H^0=0\), the de Rham filtration has
\(\operatorname{Fil}^0=0\), and Bloch--Kato together with
local Euler--Poincare duality gives
\[
 \dim_E Z^1(H,V)-\dim_E Z^1_f(H,V)
                       =\dim_E H^0(H,V^\vee(1))
 \tag{LC5}
\]
for either representation \(V\).
Coboundaries cancel in this equality:
\(\dim H^1_f=[k:\mathbb Q_\ell]\dim_E V\), whereas
\(\dim H^1=[k:\mathbb Q_\ell]\dim_E V+
                 \dim H^0(V^\vee(1))\).

If \(w\ge2\), some inertia element makes
\(\chi(g)-\epsilon_\ell(g)\) a nonzero scalar, hence a
unit.
The dual invariants vanish, and (LC5) gives \(Z^1=Z^1_f\).

Suppose \(w=1\), and put
\(\delta=\chi(\sigma)-\epsilon_\ell(\sigma)\).
The ideal generated by all
\(\chi(g)-\epsilon_\ell(g)\) is \((\delta)\).
This holds on integer powers of \(\sigma\) by the geometric
sum identity, on inertia by equality, and on their closure
by continuity in the finite-dimensional algebra.
Thus
\[
 H^0(H,B'(\chi)^\vee(1))
                         =\operatorname{Hom}_E(B'/\delta B',E),
\]
whose dimension is
\(\dim_E\operatorname{Ann}_{B'}\delta=\dim_EW\).
The action on \(W\) is precisely scalar \(\epsilon_\ell\),
so its dual-invariant dimension is also \(\dim_EW\).

The inclusion \(W\hookrightarrow B'\) induces an injection
\[
       Z^1(H,W)/Z^1_f(H,W)
             \longrightarrow Z^1(H,B'(\chi))/Z^1_f(H,B'(\chi)).
\]
To check injectivity, a cocycle on \(W\) defines the extension
of the trivial representation by \(W\).
If its pushout to \(B'(\chi)\) is crystalline, the original
extension is a subrepresentation of that crystalline
pushout, and is crystalline too.
This is exactly membership in \(Z^1_f\).
By (LC5) both quotient spaces have dimension \(\dim_EW\).
The injection is an isomorphism.
Choose a representative of the preimage of the class of \(z\);
this is \(z_{\rm exc}\), proving (LC4).
\end{proof}

\begin{srlemma}\label{lem:sr-triangular-adjacent-crystalline}
Let a continuous rank-three representation over a finite
local \(\mathbb Q_\ell\)-algebra \(B\) be
\[
 \rho(g)=
 \begin{pmatrix}
 \kappa_0(g)&x(g)\kappa_1(g)&z(g)\kappa_2(g)\\
 0&\kappa_1(g)&y(g)\kappa_2(g)\\
 0&0&\kappa_2(g)
 \end{pmatrix},
 \qquad
 \kappa_i|_{H\cap I_\ell}=\epsilon_\ell^{h_i}.
 \tag{LC6}
\]
If the two adjacent cocycles \(x,y\) belong to their finite
local subspaces, then \(\rho\) is crystalline.
\end{srlemma}
\begin{proof}
The first two coordinates give a crystalline extension
\(V_{01}\) of \(\kappa_1\) by \(\kappa_0\), by the
finite-extension interpretation of \(H^1_f\).
Set \(M=V_{01}\otimes\kappa_2^{-1}\).
It is crystalline and fits into
\[
 0\longrightarrow B(\kappa_0/\kappa_2)
       \longrightarrow M\longrightarrow B(\kappa_1/\kappa_2)
       \longrightarrow0 .
 \tag{LC7}
\]
Its positive exponents are \(2a,a\).
Matrix multiplication in (LC6) says that \((z,y)\) is
a cocycle of \(M\).

The finite cohomology map in (LC7) is surjective.
Here is the period calculation that verifies this assertion,
also for nonreduced \(B\), viewed as a finite-dimensional
coefficient-field representation.
All three crystalline modules have
\(\operatorname{Fil}^0D_{\rm dR}=0\).
On their \(D_{\rm cris}\), \(1-\varphi\) is invertible:
their Frobenius slopes are those of positive powers of
\(\epsilon_\ell\) times unramified unit characters.
The Bloch--Kato fundamental sequence therefore identifies
\(H^1_f\) functorially with \(D_{\rm dR}\).
Exactness of \(D_{\rm dR}\) proves the asserted surjectivity.
At the level of cocycles a lifted coboundary can be
subtracted to make the second coordinate equal to \(y\)
exactly.

Choose a finite cocycle of \(M\) with second coordinate \(y\).
Its difference from \((z,y)\) is a cocycle in
\(B(\kappa_0/\kappa_2)\).
The latter has exponent \(2a\ge2\); (LC5) has zero
dual-invariant term, so all its cocycles are finite.
Hence \((z,y)\) is finite.
The corresponding extension of \(\kappa_2\) by \(V_{01}\)
is crystalline, which is the assertion.
\end{proof}

\begin{srproposition}\label{prop:sr-ordinary-saturation-crystalline}
For every finite \(\mathbb Q_\ell\)-algebra \(B\) and
coefficient map \(f:\mathcal R\to B\),
\[
       I_Y\subset\ker f
       \quad\Longrightarrow\quad
       r_f|_H\text{ is crystalline}.
 \tag{LC8}
\]
This holds for every \(a\ge1\).
\end{srproposition}
\begin{proof}
A finite algebra is the product of its local Artinian
factors, and crystalline admissibility is preserved under
finite direct sums. It suffices to treat \(B\) local.
Apply Lemma~\ref{lem:sr-saturated-artinian-lift} to
\(J=I_Y\) and \(\Delta\) in (LC2).
Its saturation condition follows immediately from the definition:
\(\Delta b\in I_Y\) implies \(b\in I_Y\).
Obtain \(B',f',q\) as in (LC3).

We give the ordered local normal form over \(B'\).
Since \(J_{\rm ord}\subset I_Y\), its equations hold
at \(r_{f'}\).
The characteristic polynomial of \(r_{f'}(s_0)\) is
\(\prod_i(T-p_i)\).
The projectors
\[
                 e_i=\prod_{j\ne i}
                   \frac{r_{f'}(s_0)-p_j}{p_i-p_j}
\]
split \((B')^3\) into three rank-one summands.
They are free since \(B'\) is local.
Choose their bases and conjugate so
\(r_{f'}(s_0)=\operatorname{diag}(p_0,p_1,p_2)\).

For \(u\in H\cap I_\ell\), use the three triple products
with arguments \((u,s_0,s_0)\), \((s_0,u,s_0)\),
and \((s_0,s_0,u)\).
Their entries, after cancellation of the unit differences
\(p_i-p_j\), give
\[
 r(u)_{10}=r(u)_{20}=r(u)_{21}=0,\qquad
 r(u)_{ii}=\Psi(u)^{(1,0,-1)_i}.
 \tag{LC9}
\]
For example the first product gives the entries \(10,20,00\);
the middle product gives \(21,11\); the last gives \(22\).
For \(g\in H\), normality puts
\(gs_0g^{-1}\) back in \(H\cap I_\ell\), and its three
prescribed diagonal values are \(p_i\).
Writing \(T=r(gs_0g^{-1})\), one has
\[
                     T r(g)=r(g)\operatorname{diag}(p_i).
\]
The \(20,21,10\) entries in this order, and the same unit
differences, give those three lower entries of \(r(g)\)
equal to zero.
Thus the entire \(H\)-representation has the form (LC6).
Its diagonal entries are continuous characters
\(\kappa_i\), since multiplication of upper triangular
matrices multiplies their diagonals.
The adjacent entries \(x=r_{01}/\kappa_1\) and
\(y=r_{12}/\kappa_2\) are the corresponding cocycles.

The trace definition in (LC2) now gives the exact identities
\[
 T_i=c_i\kappa_i(\sigma),\qquad
 d_i=c_ic_{i+1}\kappa_{i+1}(\sigma)
       \left((\kappa_i/\kappa_{i+1})(\sigma)
                                     -\epsilon_\ell(\sigma)\right).
 \tag{LC10}
\]
Every factor preceding the parentheses is a unit of \(B'\).
Consequently each exceptional annihilator in
Lemma~\ref{lem:sr-local-positive-annihilator}, applied to
\(x\) and \(y\), is contained in
\(\operatorname{Ann}_{B'}f'(\Delta)\).
Its image under \(q\) vanishes by (LC3).
Equation (LC4) therefore says that \(q x\) and \(q y\)
are finite cocycles for the pushed diagonal characters.
Lemma~\ref{lem:sr-triangular-adjacent-crystalline}
makes \(q r_{f'}=r_f\) crystalline on \(H\).
Conjugating back preserves crystalline admissibility.

All representation maps used here are continuous;
we check the nilpotent target as well.
Let \(C=f(\mathcal R)\) in a finite Artinian
\(\mathbb Q_\ell\)-algebra.
In each residue-field factor every element of the
semisimple image has absolute value at most one.
Otherwise the reverse minimal polynomial of its inverse,
multiplied by a scalar to clear denominators, gives a
relation
\(1+\sum_{i\ge1}\varpi a_i r^i=0\) in that image,
with \(a_i\in\mathcal O\).
Its preimage is a unit of the local ring \(\mathcal R\),
a contradiction.
Thus the semisimple image is an \(\mathcal O\)-submodule
of the finite product of the corresponding integer
lattices, and is finite over \(\mathcal O\).

The image \(C\) is Noetherian, being a quotient of
\(\mathcal R\).
Its nilpotent radical \(\mathfrak n\) has
\(\mathfrak n^j/\mathfrak n^{j+1}\) finitely generated
over the finite \(\mathcal O\)-algebra \(C/\mathfrak n\).
There are finitely many such pieces, so \(C\) itself is
finite over \(\mathcal O\).
It is local, and its maximal-ideal and
\(\varpi\)-adic topologies are cofinal:
\(C/\varpi C\) is Artinian local, so some power of its
maximal ideal is zero.
The surjection \(\mathcal R\to C\) is local, hence
continuous for these topologies, and the finite
\(\mathcal O\)-module inclusion of \(C\) in the
coefficient algebra is continuous.
This proves the assertion for \(f\), and equally for
the thickened point \(f'\).
\end{proof}

\subsubsection{The selected component and the actual global chart}
Let \(x_0:\mathcal R\to\mathcal O\) be the torus point.
It kills \(I_Y\).
For completeness, deform \(\Psi\) by a continuous unramified
character \(\lambda_t\), with value \(1+t\) on a rational
uniformizer lift, over \(\mathcal O[[t]]\).
The torus representation for \(\Psi\lambda_t\) kills
\(J_{\rm ord}\).
Its value of \(\Delta\) is a nonzero power series:
by (LC10) it is a nonzero constant times a unit times
\[
        \left(\Psi(\sigma)(1+t)^f
                                      -\epsilon_\ell(\sigma)\right)^2
\]
for the positive residue degree \(f\) of \(k\).
Since \(\mathcal O[[t]]\) is a domain, it kills \(I_Y\)
as well.
Specialization at \(t=0\) proves the assertion about \(x_0\).
This argument includes \(a=1\), when \(\Delta(x_0)\)
itself can be zero.

Take the minimal primes of \(I_{\rm cr}\) contained in
\(\ker x_0\), and let \(I_t\) be their intersection.
The set is nonempty, since \(\ker x_0\) is prime and contains
\(I_{\rm cr}\).
Thus
\[
                       I_Y\subset I_{\rm cr}\subset I_t.
 \tag{LC11}
\]
This is exactly the source's selection \(\mathrm{primesAt0}\)
and \(\mathrm{It}\); uniqueness of the component is not
needed for (LC8) or (LC11).

We next verify the factorization at every point that is counted.
The actual global ring is not an arbitrary ordinary deformation ring.
It is the ring representing
\[
 \mathrm{NormCond}\bigl(c,g_1,
       \mathrm{BoxCond}(\mathrm{Kill},\iota,
                            \mathcal R,I_t)\bigr).
 \tag{LC12}
\]
Unwinding \(\mathrm{BoxCond}\) and \(\mathrm{OnI}\),
its universal representation carries a coefficient map
\(\alpha:\mathcal R\to R_{\rm glob}\) such that
\[
 \alpha(I_t)=0,\qquad
           \alpha_*r=r_{\rm glob}\circ\iota.
 \tag{LC13}
\]
This is part of the represented condition itself.
For every actual global point \(f:R_{\rm glob}\to B\),
the composite \(f\alpha\) therefore kills \(I_t\).
Equations (LC11) and (LC8) prove its crystalline assertion.

The same verification applies to the graded chart points
in the one-character intersection argument.
Its point comparison supplies, for every selected mixed
chart point \(\xi\), a scalar \(z\) and an actual local
point \(f:R_{\rm glob}\to B\) with
\[
                       \xi^{\,z}=f\circ q_{\rm chart}.
 \tag{LC14}
\]
This is the point comparison already used in
Section~\ref{sec:single-character};
on universal matrices the scaling is conjugation by
\(\operatorname{diag}(z,1,z^{-1})\).
Use (LC13) at \(f\) first.
Conjugating the resulting representation preserves
crystalline admissibility, so transports the conclusion
to the original chart point.
No inference is made that an unrelated ordinary point
kills \(I_t\), and no new condition is imposed after
the chart point has been chosen.

In the frozen source (LC12)--(LC14) are literally
\srcid{UnitaryThreeGroups/GRepresents.lean}
(\srcid{OnI}, \srcid{BoxCond}),
\srcid{ConjStable.lean} (\srcid{SRing}, \srcid{Model.pts}),
and \srcid{RhoKApply.lean} (\srcid{Dat.Ring}, \srcid{IsPt.loc}).
Thus the selected-chart factorization required by (GE7)
is verified on the actual rings.

\begin{srcorollary}\label{cor:sr-generic-original-level}
Assume the actual type-I automorphic endpoint of
Proposition~\ref{prop:sr-generic-three-exact-endpoint}.
Then every selected point satisfies the source's local
nonlinking assertion \(\mathrm{Th.Unl}\), also when \(a=1\).
Its automorphic representation has the original principal
congruence level \(N\) over \(K\).
In particular no varying auxiliary field or coefficient
prime is introduced into the fixed spectrum's level.
\end{srcorollary}
\begin{proof}
By (LC8)--(LC14) the local representation is crystalline
on \(H=\ker\eta\cap\ker\tau\).
Its Weil--Deligne monodromy over the original local field
is therefore zero.
The ordered flag over \(H\) descends to the original
decomposition group over \(K\).
Indeed every conjugate flag has the same ordered characters
on open inertia, namely the distinct powers
\(\epsilon_\ell^a,1,\epsilon_\ell^{-a}\);
an element with infinite cyclotomic value makes the ordered
eigenflag unique.
Thus the local parameter is a sum of three characters.

Use full local--global compatibility at \(\ell\).
The generic unitary classification gives the strict window
(GE5), now without a prior nonlinking assumption, since
\(N=0\) was proved independently.
A relation
\(\kappa_i=\kappa_j\epsilon_\ell\)
would have Frobenius absolute-value ratio \(q_v\),
contrary to that window.
This proves \(\mathrm{Th.Unl}\).
The same argument also rules out every finite-power
cyclotomic relation.

For the level assertion, if \(\ell\nmid N\), then
\(\tau=1\) and \(H\) is already the decomposition group
over \(K\).
The representation is crystalline there, with trivial
Weil--Deligne inertia; hence the automorphic component
is unramified.
If \(\ell\mid N\), its finite Weil--Deligne inertia
factors through the quadratic tame quotient cut out
by \(\psi_\ell\).
Every inducing character therefore has conductor
exponent at most one, which is at most \(v_\ell(N)\).
We do not assert unramifiedness in this case.

At \(p\ne\ell\), the original killed-inertia condition
makes inertia on \(G_K\) factor through the cyclic
character \(\bar\psi|_{I_p}\).
Its order is prime to \(\ell\), so its Teichm\"uller
lift has exactly the same kernel.
Every resulting inertial character is a power of
\(\psi|_{I_p}\), and has conductor exponent at most
\(v_p(N)\).
The full local parameter has \(N=0\) and is a sum of
characters, as in the endpoint proof.
Its generic local representation is consequently
the unlinked principal series of these characters.
A character of conductor at most \(p^{v_p(N)}\)
is trivial on the diagonal units of that principal
congruence subgroup; the usual induced-function vector
therefore gives a \(K_3(p^{v_p(N)})\)-fixed vector.
The same observation applies at \(\ell\mid N\).
At primes outside \(N\) the component is spherical.
Taking their restricted tensor product proves the
\(K_3(N)\)-fixed vector assertion.

All statements concern the descended representation over
the fixed field \(K\), at its original weight
\((a-1,-1,-a-1)\).
The auxiliary seed weight, the degree of \(L\), and the
Ferrero--Washington constant occur only in the support
construction and are absent from this level comparison.
\end{proof}

% Exact local source chain:
% ExistsFiniteLocalLiftOfSaturatedIdeal -> TriangularConjOfOrdIdealEqBot
% -> CocycleCrysAddAnn -> TriCrys -> TriangularRankThreeCrystallineAfterPush
% -> IsLCrystallinePointOfSaturatedOrdIdeal -> SaturatedOrdIdealPointsIsLCrystalline
% -> RtPoints.
